\newpage
\section{Input analysis}

\subsection{Purpose}
This section provides details for the stochastic inputs for the simulation in accordance with chapter 9 of the textbook \parencite{pensumbok}. The purpose is to identify different elements of the Miller Pain Treatment Center \parencite{MillerPainTreatmentCenter}, organize the available data, fit appropriate probability distributions and justify all assumptions made in the process.

\subsection{Inputs}
The inputs from the Miller Pain Treatment Center can be categorized into three main groups; patient arrivals, treatment times and outcomes for patients (routing decisions). Each of these groups contains multiple stochastic variables that must be modelled probabilistically to capture the random nature of the system.

Patient arrivals are stochastic as they depend on patients arriving early, on time or late for their appointments. This results in queues and system congestion, which must be captured by the model.

Treatment times also vary between patients depending on their conditions and the complexity of the treatment needed. These variables include registration time, vitals, consultation time and treatment, as well as time spent on teaching residents.

Finally, patient outcomes provides stochastic routing decisions in the system. According to their needs, patients follow different paths through the system, resulting in varying resource utilization, waiting times and overall system performance. 

All of these variables must be caught and modelled probailistically to ensure the simulation accurately depicts the system. \autoref{tab:stochastic_vars} summarizes the stochastic variables identified in the Miller Pain Treatment Center, categorized into the three main groups:

\begin{table}[H]
\centering
\small
\begin{tabularx}{\textwidth}{p{4cm} p{2cm} X}
\toprule
\textbf{Variable} & \textbf{Type} & \textbf{Description} \\
\midrule

\multicolumn{3}{l}{\textbf{Service times}} \\
Registration time & Continuous & Time spent on administrative registration \\
Vitals time & Continuous & Time to record patient vitals \\
Physician assistant time & Continuous & Time spent by PA on follow-up patients \\
Check-out time & Continuous & Time to complete visit and paperwork \\
Resident review (new) & Continuous & Time reviewing patient file (new patients) \\
Resident--patient interaction (new) & Continuous & Consultation time with resident \\
Teaching time (new) & Continuous & Teaching interaction between resident and attending \\
Resident--attending interaction (new) & Continuous & Discussion between resident and attending \\
Resident review (return) & Continuous & Time reviewing file (return patients) \\
Resident--patient interaction (return) & Continuous & Consultation time with resident \\
Teaching time (return) & Continuous & Teaching interaction for return patients \\
Resident--attending interaction (return) & Continuous & Discussion between resident and attending \\

\midrule
\multicolumn{3}{l}{\textbf{Routing decisions}} \\
Patient type & Discrete & New, return, or follow-up patients \\
Follow-up requires attending & Discrete & Whether PA consults attending (approximately 50\%) \\
No-show / cancellation & Discrete & Whether patient attends appointment (approximately 10\% no-show) \\

\bottomrule
\end{tabularx}
\caption{Identified stochastic inputs}
\label{tab:stochastic_vars}
\end{table}